\documentclass[11pt,a4paper]{article}
\usepackage[T1]{fontenc}
\usepackage[utf8]{inputenc}
\usepackage{lmodern}
\usepackage[english]{babel}
\usepackage{amsmath,amssymb,geometry,enumitem}
\geometry{margin=24mm}
\setlength{\parindent}{0pt}
\setlength{\parskip}{6pt}
\newcommand{\match}{\(\checkmark\)\;Match.}
\title{Linear Optimization (311.160)\\Exercise sheet 1}
\author{Raphael Bleier}
\date{6 October 2026}
\begin{document}
\maketitle
\section*{Exercise 1.1: definitions and comparison}
The AI answers below use the maximization convention from the current lecture notes. For every item, the comparison refers to \emph{Lecture 2 -- 6.10.2026}, PDF pages 5, 8 and 9. Each answer is followed by the corresponding lecture definition and a check.
\begin{enumerate}[leftmargin=*,itemsep=10pt]
\item \textbf{Standard form.}

\textit{AI answer:} A linear program in standard form is
\[
\max\;c^{\mathsf T}x\quad\text{subject to}\quad Ax=b,\quad x\geq0,
\]
where $A\in\mathbb R^{m\times n}$, $b\in\mathbb R^m$ and $c,x\in\mathbb R^n$. The objective is linear, the constraints $Ax=b$ are equalities, and all variables are nonnegative.

\textit{Lecture comparison (p.~5):} The notes give precisely this maximization problem as standard form. \match\ Other textbooks may use a different convention, but that is not the convention used here.
\item \textbf{Feasible or infeasible vector.}

\textit{AI answer:} A vector $x\in\mathbb R^n$ is feasible if it satisfies every constraint, that is, $Ax=b$ and $x\geq0$. It is infeasible if at least one constraint fails.

\textit{Lecture comparison (p.~8):} The notes call a vector satisfying all constraints feasible and any other vector infeasible. \match
\item \textbf{Feasible set.}

\textit{AI answer:} The feasible set contains all feasible vectors:
\[
C=\{x\in\mathbb R^n:Ax=b,\ x\geq0\}.
\]
\textit{Lecture comparison (p.~8):} The notes define the same set $C$. \match
\item \textbf{Feasible or infeasible linear program.}

\textit{AI answer:} The LP is feasible if $C\neq\varnothing$, meaning that some vector satisfies all constraints. It is infeasible if $C=\varnothing$.

\textit{Lecture comparison (p.~8):} The distinction is exactly $C\neq\varnothing$ versus $C=\varnothing$. \match
\item \textbf{Optimal value.}

\textit{AI answer:} For a maximization LP, the optimal value is
\[
\gamma=\sup\{c^{\mathsf T}x:x\in C\}\in\mathbb R\cup\{-\infty,+\infty\},
\qquad \sup\varnothing=-\infty.
\]
Thus an infeasible LP has value $-\infty$, a feasible unbounded LP has value $+\infty$, and a feasible bounded LP has a finite real optimal value.

\textit{Lecture comparison (pp.~8--9):} The notes use this supremum and the same empty-set convention, then list these three cases. \match\ The supremum definition does not itself assert that a maximizer exists.
\item \textbf{Bounded or unbounded linear program.}

\textit{AI answer:} A \emph{feasible} maximization LP is unbounded if, for every $M\in\mathbb R$, there is an $x\in C$ with $c^{\mathsf T}x>M$. It is bounded if some $M\in\mathbb R$ satisfies $c^{\mathsf T}x\leq M$ for every $x\in C$.

\textit{Lecture comparison (p.~8):} These are exactly the two conditions in the notes. \match\ Both definitions assume feasibility; an infeasible LP is a separate case. Boundedness here concerns the objective, not whether the feasible set is a bounded set.
\item \textbf{Optimal solution.}

\textit{AI answer:} A vector $x^*\in C$ is an optimal solution, or maximizer, if
\[
c^{\mathsf T}x^*\geq c^{\mathsf T}x\quad\text{for every }x\in C.
\]
Equivalently, it is feasible and attains the optimal value.

\textit{Lecture comparison (p.~9):} The notes use the same comparison with every feasible vector and also call such a vector a maximizer. \match
\end{enumerate}
\newpage
\section*{Exercise 1.2: choosing information sources}
Let $x_i\in\{0,1\}$ indicate whether option $i$ is chosen:
\begin{center}
\begin{tabular}{clrr}
$i$ & Option & Value & Time (min)\\\hline
1 & Ask a colleague & 4 & 5\\
2 & Attend the first lecture & 8 & 9\\
3 & Read the Moodle course description & 7 & 7\\
4 & Read the LV-Karte & 16 & 15\\
5 & Do nothing & 40 & 30
\end{tabular}
\end{center}
Each option is chosen at most once. Values and times add across chosen options.

\textbf{Main interpretation: at least two actual information sources.}
Doing nothing provides no information and therefore does not count as a source. The requirement to use at least two sources is mandatory, not waived by option 5. The binary optimization model is
\begin{align*}
\max\quad &4x_1+8x_2+7x_3+16x_4+40x_5\\
\text{subject to}\quad
&5x_1+9x_2+7x_3+15x_4+30x_5\leq30,\\
&x_1+x_2+x_3+x_4\geq2,\\
&x_i\in\{0,1\},\qquad i=1,\ldots,5.
\end{align*}
Option 5 already consumes the entire time budget. Every actual source takes positive time, so $x_5=1$ contradicts the requirement to use two sources. Hence every feasible choice has $x_5=0$; retaining it in the model makes that consequence explicit. If all five listed options are instead counted toward the condition $\sum_{i=1}^5x_i\geq2$, the feasible choices are unchanged because option 5 still cannot be combined with another option.

As a check, the best feasible choice is $(x_1,x_2,x_3,x_4,x_5)=(1,1,0,1,0)$: it uses three information sources, takes $5+9+15=29$ minutes, and has value $4+8+16=28$.

\textbf{Alternative interpretation: doing nothing is an opt-out.}
The inclusion of option 5 could mean that the student may decline to gather information altogether. That reading needs an explicit exception to the two-source requirement. Under that assumption, replace the source constraint by
\[
x_1+x_2+x_3+x_4\geq2(1-x_5),
\qquad x_i+x_5\leq1\quad(i=1,\ldots,4),
\]
while keeping the objective, time budget and binary domains. Then either at least two sources are used, or only option 5 is chosen. The latter gives value 40 in 30 minutes and is optimal for this alternative model. It does \emph{not} satisfy the sheet's literal two-source requirement, so the main model above is my formulation without an opt-out exception.

This is a binary integer linear optimization problem, following the selection model in Lecture 1, p.~8. No conversion to standard form is needed.
\newpage
\section*{Exercise 1.3: FreshShine}
Define the monthly production quantities
\begin{align*}
x_1&=\text{units of Citrus Shower Gel},\\
x_2&=\text{units of Citrus Shampoo},\\
x_3&=\text{units of Citrus Body Lotion}.
\end{align*}
For the LP formulation, quantities are divisible and $x_1,x_2,x_3\in\mathbb R_{\geq0}$. If whole units were required, integer domains would be an additional assumption.

Since all produced units can be sold, the profit in euros is $3x_1+5x_2+4x_3$. The model is
\begin{align*}
\max\quad &3x_1+5x_2+4x_3\\
\text{subject to}\quad
&x_1+2x_2+x_3\leq170 &&\text{(machine A, hours/month)},\\
&x_1+x_2+2x_3\leq150 &&\text{(machine B, hours/month)},\\
&3x_2+x_3\leq180 &&\text{(machine C, hours/month)},\\
&x_1,x_2,x_3\geq0.
\end{align*}
Citrus Shower Gel uses no machine C time, which explains its zero coefficient in that constraint. There is no sales constraint because the sheet states that every produced unit can be sold. This follows the production-model structure in Lecture 1, pp.~3--4: nonnegative quantities, total profit, and one constraint for each limited resource. No conversion to standard form is needed.

\end{document}
